% Variante de non-regression : numerotation dans un TD sans section
% (ocots-latex-template#67). Lue par check-td-numbering.sh :
%   - les boites ne portent pas de numero de section (« Remarque 1 », pas
%     « Remarque 0.1 ») ;
%   - une question et une sous-question se citent par \ref, avec le numero
%     affiche (2.2, 2.2.2).
\documentclass[11pt]{ocots-td}
\usepackage[lang=fr, theme=ocots, solutions=inline, math=analysis, institution={n7}]{ocots}

\title{Variante : numérotation dans un TD}
\shorttitle{Variante : numérotation dans un TD}
\numero{TD 0}
\date{\today}
\discipline{Template ocots}
\promotion{Harnais de non-régression}

\begin{document}

\maketitle

\begin{exercise}[label=ex:un]
    Premier exercice.
    \begin{question}
        Une question.
    \end{question}
\end{exercise}

\begin{exercise}[label=ex:deux]
    Second exercice.
    \begin{question}
        Première question.
    \end{question}
    \begin{question}\label{q:deux-deux}
        Seconde question.
        \begin{subquestion}\label{sq:un}
            Première sous-question.
        \end{subquestion}
        \begin{subquestion}\label{sq:deux}
            Seconde sous-question.
        \end{subquestion}
    \end{question}
    \begin{correction}
        La question~\ref{q:deux-deux} se cite depuis le corrigé.
    \end{correction}
\end{exercise}

\begin{theorem}[label=thm:td]
    Un résultat.
\end{theorem}

\begin{example}[label=exa:td]
    Un exemple.
\end{example}

\begin{remark}[label=rem:td]
    Une remarque.
\end{remark}

NUMCHECK Exercice = \ref{ex:deux}

NUMCHECK Question = \ref{q:deux-deux}

NUMCHECK Sous-question = \ref{sq:deux}

NUMCHECK Théorème = \ref{thm:td}

NUMCHECK Exemple = \ref{exa:td}

NUMCHECK Remarque = \ref{rem:td}

\end{document}
